\begin{otherlanguage*}{english}
\noindent\textbf{Accessible conditions, categories of structures, and gluing.}

These pages turn accessibility into a working calculus. Their central question is how a condition on the objects of a large category can be recognized, transported, and imposed without losing a small description of the resulting category. The first part completes a criterion for accessible full subcategories. The second applies it to functors preserving specified cones and cocones, to categories of sections, and finally to categories lying over a small base. The progression is deliberate: once membership in a subcategory can be read on a suitable presentation, many apparently new constructions become instances of the same argument.

One convention must be fixed before comparing the manuscript with current accounts. Grothendieck calls an indexing category \emph{grande devant $\pi$} when it admits cocones on diagrams having at most $\pi$ arrows. Thus the associated regular cardinal in present terminology is $\lambda=\pi^+$, and his $\pi$-accessible objects are $\lambda$-presentable objects. A useful $\pi$ provides an essentially small full subcategory $\mathcal M_\pi$ generating $\mathcal M$ under these filtered colimits. The stronger phrase accessible \emph{in} an ambient category keeps track of the inclusion and of ambient presentability, not just an abstract equivalence class of accessible categories.

The opening argument is a transfinite alternation between two presentations. A final functor $\varphi:I\to J$ and a full replete final subcategory $J'\subset J$ do not in general have a final strict inverse image: the even and odd natural numbers in their common ordered set already show the problem. Grothendieck adds an ordinal-colimit hypothesis. Starting from a chosen object of $I$, he constructs alternating maps
\[
 \varphi(i_\beta)\longrightarrow j_\beta
 \longrightarrow\varphi(i_{\beta+1}),
 \qquad j_\beta\in J',
\]
whose composites are transition maps. At the final limit, the two chains have isomorphic colimits. One colimit comes from $I$, and the other lies in $J'$; repleteness then gives an object of the inverse image beyond the prescribed starting object. Lemma 4.13.17 is therefore a genuine cofinality statement with an additional closure hypothesis. The commentary supplies its final comparison explicitly and shows that the successor-stage equalizations in the manuscript are redundant.

This lemma drives Theorem 4.13.22. Suppose $\underline S$ is an accessible full subcategory of $\mathcal M$ closed under the indicated colimit shapes. In an adapted presentation $x=\varinjlim_i x_i$, membership $x\in\underline S$ is equivalent to finality of the full subcategory of indices for which $x_i\in\underline S$. Adaptedness says that the indexing category contains the specified colimits and that the presenting diagram preserves them. It is what prevents a presentation from hiding the condition at every finite stage. The worked example of divisible abelian groups exhibits the criterion through countable divisible subgroups; the familiar presentation of $\mathbb Q$ by cyclic groups explains why an arbitrary filtered presentation is insufficient.

The subsequent calculus treats enlargement of the accessibility rank, small unions and intersections, and inverse images under accessible functors. It also exposes a boundary of the method. The essential full image of an accessible functor need not be accessibly embedded, even after closing under retracts. We work out the split-monomorphism example recorded in \citet{Rosicky2009Combinatorial,BrookeTaylorRosicky2017}: forgetting chosen retractions allows diagrams whose colimit no longer splits. This failure is distinct from the controlled envelope construction. Starting with a small family of objects and closing it under a fixed filteredness threshold and a small set of additional colimit shapes does yield an accessible, accessibly embedded subcategory. The commentary proves this by choosing a sufficiently closed cardinal, making bounded filtered envelopes of index categories, and identifying the result with a fully faithful Ind-image.

The historical framework is already visible in \citet[Expos\'e I, 8.7]{SGA4I1972}, where Ind-objects and their universal extension property are developed. The systematic accessible-category theory of \citet{MakkaiPare1989} was published in 1989, before these pages were written. The relation is therefore a comparison of formulations, not a claim that this manuscript first announced a theory later discovered by those authors. The modern presentation separates products, inserters, equifiers, and pseudopullbacks; Grothendieck often constructs the same operations through accessible loci in categories of diagrams. \citet{Positselski2024Limits} makes especially clear how additional chain-colimit hypotheses permit control of a specified rank.

Section 4.14 begins with marked cones and cocones in a small category. A functor is required to take the marked data to genuine limits or colimits. When the target has the relevant universal constructions, each condition is the inverse image of an isomorphism locus under an accessible comparison functor. Their small intersection is again accessible. This is the language of sketches and categories of models. Restriction to presentable objects then allows certain large-source functor categories to be treated by the same method.

Here the hypotheses require particular care. The proof of the fixed-threshold functor-category assertion invokes target colimits not supplied by accessibility alone. We distinguish the theorem established with the additional target-colimit hypothesis from the more general assertion the manuscript leaves to its later question about colimit cones. We also distinguish the two adjoint theorems: cocontinuous functors from a locally presentable category have right adjoints, whereas a limit-preserving covariant functor need not have a left adjoint unless a solution-set or accessibility condition is added. The stronger claim at pages 379--380 fails already for the category of groups, a warning that Grothendieck himself recalled near the beginning of this chapter. The counterexample and the appropriate repair are explained with the primary references.

The question about epimorphisms has a different status. Grothendieck observes that epimorphisms are stable under the relevant filtered colimits and asks whether their full subcategory in the arrow category is always accessible. Closure alone does not answer that question. The commentary proves that accessibility and accessible embedding of this epimorphism locus would force the underlying accessible category to be co-wellpowered. Published results show that the universal co-wellpoweredness assertion already has a large-cardinal lower bound, while stronger large-cardinal assumptions give positive results \citep{BagariaCasacubertaMathiasRosicky2015}. This identifies a real set-theoretic obstruction. It does not settle the exact strength of the epimorphism-locus assertion. The recent study of pure epimorphisms by \citet{Positselski2026Pure} concerns an additional purity condition; its publisher abstract was checked as a status lead, and no result about the unrestricted epimorphism class is inferred from it.

Categories of sections provide the main constructive application. For a fibration, a prospective section consists of a family of fibre objects and arrows $X_s\to\lambda^*X_t$. Those data form an inserter. Identity and composition impose a small family of equations, and cartesian or cocartesian requirements impose isomorphism conditions. The proof keeps track of a nonsplit cleavage through its canonical coherence isomorphisms. It is thereby a genuinely categorical construction, not a pointwise prescription with the functoriality omitted. The higher-category analogue is the accessible section theorem of \citet[Proposition 5.4.7.11]{Lurie2009HTT}.

The qualitative section theorem is robust; several printed rank assertions are not. For unrestricted sections of an opfibration, pointwise presentability need not imply presentability of the section unless transport has an appropriate presentability bound. Conversely, the category of sections can be very small while a forced component is large. The commentary gives explicit examples for both directions. Cocartesian sections admit a separate argument, because their structure isomorphisms identify the transported source with the target component. Keeping that distinction avoids discarding a correct corollary along with an incorrect proof reference.

The final pages replace fibrations by intervalwise accessible data. An arbitrary functor $p:\mathcal E\to\mathcal C$ determines, over each base arrow, a category obtained by gluing its two endpoint fibres. Its cross-Hom bifunctor
\[
 H:\mathcal E_0^{\mathrm{op}}\times\mathcal E_1\to\mathrm{Set}
\]
controls the gluing. Contravariant preservation makes fibrewise colimits universal in the total category; covariant preservation makes source objects presentable at a suitable rank. This proves the accessible-section theorem without requiring chosen lifts. An additional hypothesis on idempotents in the base addresses accessibility of the total category itself.

Two corrections matter for reusing this last argument. Corollary 4.14.25 omits the colimit hypothesis on the closed fibre from one side of its equivalence; an ordinal example disproves the printed version. Corollary 4.15.2 also fails to put the contravariant preservation ranks below the initial presentation rank. The power-functor collage $X\mapsto X^{\aleph_1}$ gives a counterexample satisfying the printed cardinal bounds. Both qualitative constructions survive a common-rank choice with the missing requirements included. These pages end during the proof; the complete chapter later justifies the constant-subdiagram argument with the strict inequality $\pi>|\mathcal C|$. That later repair is reported at its exact manuscript locations.

The audit records these distinctions assertion by assertion: what the source proves, sketches, states, or asks; what is merely a slip in the edition; and what needs a mathematical repair. A few exact formulations remain uncertified, including the unrestricted fixed-rank third variant of 4.13.21 and the broad weakening suggested in footnote 172. These are identified as limits of the present verification rather than declared counterexamples or new open problems. Proofs marked \emph{[ours]} are arguments supplied by the commentators, with the scope of the literature check stated and no claim of priority.

\medskip
\noindent\textit{Credit.} Commentary, audit, and mathematical developments AI-produced by GPT-6 Astra (OpenAI, ChatGPT Pro mode), checked and integrated by Claude (Anthropic).
\end{otherlanguage*}
